\documentclass[11pt,a4paper]{article}
\usepackage{fontspec,amsmath,amssymb,amsthm,microtype}
\setmainfont{lmroman10-regular.otf}[BoldFont=lmroman10-bold.otf,ItalicFont=lmroman10-italic.otf,BoldItalicFont=lmroman10-bolditalic.otf]
\usepackage{xcolor}
\usepackage[most]{tcolorbox}
\usepackage{mathrsfs}
\providecommand{\lt}{<}
\providecommand{\gt}{>}
\setmonofont{DejaVuSansMono.ttf}[BoldFont=DejaVuSansMono-Bold.ttf,ItalicFont=DejaVuSansMono-Oblique.ttf,BoldItalicFont=DejaVuSansMono-BoldOblique.ttf,Scale=0.86]
\newfontfamily\leanfallback{STIXTwoMath-Regular.otf}[Scale=0.86]
\definecolor{leanrule}{gray}{0.62}
\definecolor{leanback}{gray}{0.962}
\definecolor{leantitle}{gray}{0.22}
\definecolor{leanmuted}{gray}{0.38}
\newtcolorbox{leanbox}[1]{enhanced,breakable,colback=leanback,colframe=leanrule,
  boxrule=0.4pt,arc=0.6mm,left=2.4mm,right=2.4mm,top=1.8mm,bottom=1.8mm,
  before skip=0.8em,after skip=1em,fonttitle=\small\bfseries,coltitle=leantitle,
  colbacktitle=leanback,titlerule=0.3pt,title={#1}}
\newlength{\leancharwidth}
\newenvironment{leancode}{\par\smallskip\begingroup\ttfamily\footnotesize\raggedright
  \settowidth{\leancharwidth}{M}\setlength{\parskip}{0pt}\setlength{\parindent}{0pt}
  \frenchspacing}{\par\endgroup\smallskip}
\newcommand{\leanline}[2]{\hangindent=\dimexpr#1\leancharwidth+4\leancharwidth\relax
  \hangafter=1\noindent\hspace*{#1\leancharwidth}#2\par}
\newcommand{\leansource}[1]{\par{\footnotesize\color{leanmuted}#1}\par}
\newcommand{\leanlabel}[1]{\par\smallskip\noindent{\small\bfseries #1}\par}
\usepackage[margin=25mm]{geometry}
\usepackage[unicode,hidelinks]{hyperref}
\hypersetup{pdfauthor={LangevinLean Project}}
\setlength{\emergencystretch}{3em}
\setlength{\parskip}{0.4em}
\setlength{\parindent}{0pt}
\allowdisplaybreaks[2]
\raggedbottom
\title{\texorpdfstring{Kinetic interacting particle Langevin Monte Carlo: the results of this report}{Kinetic interacting particle Langevin Monte Carlo: the results of this report}}
\author{LangevinLean Project}
\date{}
\begin{document}\maketitle
Draft mathematical exposition; these notes may contain mistakes.


This report treats the kinetic interacting particle Langevin Monte Carlo methods of Valsecchi Oliva and Akyildiz for maximum marginal likelihood estimation. A joint potential \(U(\theta,x)=-\log p_\theta(x,y)\) on \(\mathbb R^{d_\theta}\times\mathbb R^{d_x}\) is given; the goal is the maximizer \(\theta_\star\) of the marginal likelihood \(\ell(\theta)=\int e^{-U(\theta,x)}\,dx\). The standing assumptions are that \(U\) is differentiable and \(m\)-strongly convex with \(m>0\), that \(\nabla U\) is \(L\)-Lipschitz, and that the number of particles \(N\) is positive. The particle state is written in the scaled coordinates \(z=(\theta,X^1/\sqrt N,\dots,X^N/\sqrt N)\) of dimension \(d_\theta+Nd_x\), in which the force of the kinetic particle diffusion is the gradient of the averaged potential \(\bar U_N\), an \(m\)-strongly convex function with \(L\)-Lipschitz gradient.




The main results are the following. (1) The phase law \(\mathcal N(0,N^{-1}I)\otimes\pi_N\), with \(\pi_N\propto e^{-N\bar U_N}\), is stationary for the kinetic particle diffusion, and its parameter marginal \(\pi_\Theta\) has density proportional to \(\ell(\theta)^N\) (Proposition 1 of the paper). (2) \(\int\|\theta-\theta_\star\|^2\,d\pi_\Theta\le d_\theta/(Nm)\) (Proposition 2). (3) If \(\gamma^2\ge m+L\), the diffusion started with independent thermal velocities converges to \(\pi_N\) in Wasserstein distance at rate \(\sqrt2e^{-mt/\gamma}\), and its parameter marginal satisfies the bound of Proposition 3. (4) For KIPLMC1, under \(\gamma^2\ge m+L\) and \(4\gamma Lh\le m\), the root mean square error of \(\theta_n\) about \(\theta_\star\) is at most \(\sqrt2(1-\frac{3mh}{4\gamma})^nW_2(\nu,\pi_N)+\frac{\sqrt2Lh}m\sqrt{d_\theta/N+d_x}+\sqrt{d_\theta/(Nm)}\) (Theorem 1). (5) For KIPLMC2, under \(\gamma^2\ge4L\) and \(33\gamma^3h\le m\), an analogous bound holds with contraction factor \((1-\frac{mh}{3\gamma})^{n/2}\) and bias of order \(h\sqrt{d_\theta/N+d_x}\); it implies the bound of Theorem 2 of the paper with an explicit constant depending only on \(m\), \(L\), \(\gamma\). (6) Explicit choices of \(N\), \(h\) and \(n\) reach a prescribed accuracy \(\varepsilon\). (7) The Bayesian logistic regression model of the experiments has a Lipschitz gradient and is strictly convex when the covariates span the latent space, but it is not strongly convex without additional regularization.




The sections are ordered as follows: the particle geometry and the two update maps; the Gibbs phase target and the concentration of its parameter marginal; integral solutions and Markov laws; the correlated innovation measure and the law of the diffusion; the clock change that identifies both algorithms with standard kinetic algorithms for the sum potential; the error bounds for KIPLMC1 and KIPLMC2; the Brownian particle diffusion and its stationarity; the accuracy budgets; and the logistic example.


\end{document}
